
\subsection{数值与物理检验、规划--执行偏差诊断}
为验证模型求解结果满足物理约束与数值可靠性，对Q1--Q4四个问题的仿真输出逐项校验。检验维度包含：母线功率平衡残差、蓄电池SOC时序递推校验、变量上下界约束、储能充放电互斥约束、储能跨日状态连续传递、全周期逐时段现金流水与汇总核算一致性。
同时开展\textbf{规划--执行偏差诊断}：对比MPC规划输出的储能状态与真实工况执行后的SOC，计算一步预测误差。该指标用于佐证仿真流程的因果时序约束：若存在未来信息泄露，规划值将高度贴近真实执行值，一步误差会异常偏小。

对全量求解结果校验可得：全部方案的储能SOC时序递推、母线功率平衡残差均处于$10^{-13}$量级，属于求解器浮点计算误差范围；全部方案满足储能充放电互斥约束$\min(x_t,y_t)=0$，不存在同一时隙同时充放电的数值伪解。仿真全程储能SOC始终处于$1200\le S_t\le10800\ \mathrm{kWh}$的硬件边界之内。
各方案规划--执行一步SOC误差如表\ref{tab:exec_diag}所示，一步MAE为$5$--$6\ \mathrm{kWh}$量级，RMSE为$17$--$20\ \mathrm{kWh}$量级，存在合理的执行偏差，符合"决策时刻仅使用历史与已发布预测、真实值仅事后用于执行修正"的设定，未出现利用未来真值的异常现象。
\begin{table}[htbp]
\centering
\caption{各方案规划--执行一步SOC偏差诊断}
\label{tab:exec_diag}
\begin{tabular}{lccc}
\toprule
方案 & 一步MAE/kWh & 一步RMSE/kWh & 最大绝对误差/kWh \\
\midrule
Q2主策略 & 5.327 & 19.666 & 284.583 \\
Q3 $\mathcal{S}=\{6,12,18\}$ & 5.062 & 17.627 & 271.979 \\
Q4-2 固定合同 & 4.796 & 17.303 & 271.979 \\
Q4-3 可调合同 & 5.157 & 18.068 & 271.979 \\
\bottomrule
\end{tabular}
\end{table}
全年度逐时段现金流累加结果与年度汇总总费用的最大绝对偏差为$0$元，费用分项核算完全一致，不存在口径冲突。综合以上校验，各问题调度输出满足全部物理约束，数值求解误差很小，规划与执行之间存在合理偏差，仿真结果具备可信性。

\subsection{尾部风险与跨月稳定性检验}
仅依靠年度平均总费用不足以完整评价调度策略的经济性能，进一步对日运行费用分布做尾部风险度量，采用95\%分位VaR、CVaR以及最坏单日费用评估极端工况下的风险水平，结果如表\ref{tab:tail_risk}所示。
\begin{table}[htbp]
\centering
\caption{各方案尾部风险指标（95\%置信水平）}
\label{tab:tail_risk}
\begin{tabular}{lccc}
\toprule
方案 & $\mathrm{VaR}_{0.95}$/元 & $\mathrm{CVaR}_{0.95}$/元 & 最坏单日费用/元 \\
\midrule
Q2主策略 & 66075.35 & 70048.64 & 88726.19 \\
Q3 $\mathcal{S}=\{6,12,18\}$ & 65177.30 & 70474.34 & 77111.50 \\
Q4-2 固定合同 & 87294.40 & 101705.29 & 134220.63 \\
Q4-3 可调合同 & 76237.79 & 83469.29 & 92033.53 \\
\bottomrule
\end{tabular}
\end{table}
对比Q4-2与Q4-3可看出，启用日内合同调整后，$\mathrm{CVaR}_{0.95}$由101705.29元降至83469.29元，下降18235.99元；最坏单日费用由134220.63元降至92033.53元。合同调整机制不仅降低平均费用，也显著压缩极端高费用日风险，策略具备良好抗风险能力。

进一步开展跨月稳定性检验，考察各月份启用全部日内预报相对不使用新预报（Q3空集）的费用改善幅度。2--12月各月日均费用差均为负值，即每月均实现成本节约，收益覆盖全年所有月份，并非由个别异常月份单独驱动，策略改善效果具备时间维度的稳健性。

\subsection{结构敏感性分析}
为评估模型输出结果对关键参数的鲁棒性，开展多维度敏感性压力测试。\textbf{重要测试原则}：全部敏感性实验均在相同负荷序列、电价信息集、结算规则下完整重跑全周期仿真，不采用"直接替换参数代入原有调度轨迹"的简化估算，保证各组对比公平。测试变量包含：充放电效率、仿真初始SOC、储能容量/功率上限、负荷预测误差倍率，结果如图\ref{fig:sensitivity_all}所示。

\begin{figure}[htbp]
\centering
\includegraphics[width=0.85\textwidth]{figures/sensitivity-overview.png}
\caption{多维度参数敏感性测试结果汇总}
\label{fig:sensitivity_all}
\end{figure}

\subsubsection{充放电效率敏感性}
对比两组充放电效率配置$(\eta_\mathrm{ch},\eta_\mathrm{dis})=(0.9,0.9)$与$(\sqrt{0.9},\sqrt{0.9})$。效率取$\sqrt{0.9}$时，年度总费用由基准的14786101.35元降至14322473.62元，下降463627.73元（$-3.14\%$）。效率提高减少储能充放电过程损耗，系统运行费用随之下降，充放电效率属于中等敏感参数。

\subsubsection{初始SOC敏感性}
分别取储能初始SOC为额定容量的20\%、50\%、80\%。三种初始状态下年度总费用均为14786101.35元，应急购电量也完全一致，说明经过1月预热周期后，初始SOC差异对正式统计窗口的影响被完全消除，全年调度结果对初始荷电状态不敏感。

\subsubsection{储能容量与功率上限敏感性}
将储能容量、充放电功率上限分别取基准值的80\%、100\%、120\%。容量降低20\%使总费用升至15392912.44元（$+4.10\%$），容量增加20\%降至14267282.70元（$-3.51\%$）。相比之下，\textbf{功率上限降低20\%}使总费用大幅升至16318749.90元（$+10.37\%$），应急购电量由245483.29增至563286.32 kWh；功率增加20\%仅带来$-0.33\%$的微幅改善。储能功率上限属于高敏感参数，功率不足会显著限制储能移峰能力，推高应急购电成本；容量参数敏感度相对更低。

\subsubsection{负荷预测误差倍率压力测试}
设置负荷预测误差放大倍率$\gamma\in\{0.5,1.0,1.5,2.0\}$。随误差倍率增大，年度总费用呈现先升后波动：$\gamma=0.5$时费用为16632993.38元（$+12.49\%$），$\gamma=1.0$为基准值，$\gamma=1.5$时14963927.62元（$+1.20\%$），$\gamma=2.0$时15453352.16元（$+4.51\%$）。预测误差的放大或缩小都会改变风险裕度的匹配度，费用在基准倍率附近取得较低水平。在$\gamma\le2.0$范围内，本滚动风险修正框架仍能维持良好的经济性能，未出现求解失效。

综合全部敏感性实验结果：
\begin{enumerate}
    \item 储能充放电功率上限、负荷预测误差倍率属于高敏感参数；充放电效率为中等敏感参数；全年仿真结果对初始SOC敏感度最低，对储能额定容量敏感度相对较弱。
    \item 1月预热机制能有效消除仿真初始SOC带来的短期扰动，长期年度结论不受初始储能状态干扰。
    \item 在合理预测扰动范围内，整套滚动风险修正与合同调整框架可稳定发挥能量转移与风险对冲作用，同时上述结果明确了模型策略的适用边界。
\end{enumerate}
